\subsection*{Rings}
\label{subsec:rings}

\begin{definitionbox}{Definition (Ring)}
\begin{definition}[Ring]
\label{def:algebraic-ring}
A \textbf{ring} is a triple $(R,+,\cdot)$ where $R$ is a set and $+$ and
$\cdot$ are binary operations on $R$ satisfying:
\begin{enumerate}[label=\textbf{R\arabic*.}]
  \item $(R,+)$ is an abelian group;
  \item multiplication is associative;
  \item multiplication distributes over addition on the left and on the right.
\end{enumerate}
\end{definition}
\end{definitionbox}

\begin{remark*}[Standard quantified statement]
\[
\operatorname{Ring}(R,+,\cdot)
\Longleftrightarrow
\operatorname{AbelianGroup}(R,+)
\land
\operatorname{Associative}(\cdot,R)
\land
\operatorname{LeftDistributive}(\cdot,+,R)
\land
\operatorname{RightDistributive}(\cdot,+,R).
\]
\end{remark*}

\begin{remark*}[Definition predicate reading]
$(R,+,\cdot)$ is a ring exactly when addition forms an abelian group,
multiplication is associative, and multiplication distributes over addition.
\end{remark*}

\begin{remark*}[Negated quantified statement]
\[
\neg\operatorname{Ring}(R,+,\cdot)
\Longleftrightarrow
\neg\operatorname{AbelianGroup}(R,+)
\lor
\neg\operatorname{Associative}(\cdot,R)
\lor
\neg\operatorname{LeftDistributive}(\cdot,+,R)
\lor
\neg\operatorname{RightDistributive}(\cdot,+,R).
\]
\end{remark*}

\begin{remark*}[Interpretation]
A ring is the algebraic setting where addition, subtraction, and multiplication
coexist coherently.
\end{remark*}

\begin{dependencies}
\begin{itemize}
  \item \hyperref[def:algebraic-abelian-group]{Abelian Group}
  \item \hyperref[def:algebraic-associativity]{Associativity}
  \item \hyperref[def:algebraic-left-distributivity]{Left Distributivity}
  \item \hyperref[def:algebraic-right-distributivity]{Right Distributivity}
\end{itemize}
\end{dependencies}

\begin{definition}[Commutative Ring]
\label{def:algebraic-commutative-ring}
A ring $(R,+,\cdot)$ is \textbf{commutative} if
\[
\forall a,b \in R,\quad a\cdot b=b\cdot a.
\]
\end{definition}

\begin{remark*}[Standard quantified statement]
\[
\operatorname{CommutativeRing}(R,+,\cdot)
\Longleftrightarrow
\operatorname{Ring}(R,+,\cdot)
\land
\operatorname{Commutative}(\cdot,R).
\]
\end{remark*}

\begin{remark*}[Definition predicate reading]
A ring is commutative exactly when its multiplication commutes.
\end{remark*}

\begin{remark*}[Negated quantified statement]
\[
\neg\operatorname{CommutativeRing}(R,+,\cdot)
\Longleftrightarrow
\neg\operatorname{Ring}(R,+,\cdot)
\lor
\exists a,b\in R,\; a\cdot b\neq b\cdot a.
\]
\end{remark*}

\begin{remark*}[Interpretation]
Commutative rings are rings whose multiplication behaves symmetrically in its
two inputs.
\end{remark*}

\begin{dependencies}
\begin{itemize}
  \item \hyperref[def:algebraic-ring]{Ring}
  \item \hyperref[def:algebraic-commutativity]{Commutativity}
\end{itemize}
\end{dependencies}

\begin{definition}[Ring with Unity]
\label{def:algebraic-ring-with-unity}
A ring $(R,+,\cdot)$ is a \textbf{ring with unity} if there exists
$1\in R$ such that
\[
\forall a \in R,\quad 1\cdot a=a\cdot 1=a.
\]
\end{definition}

\begin{remark*}[Standard quantified statement]
\[
\operatorname{RingWithUnity}(R,+,\cdot,1)
\Longleftrightarrow
\operatorname{Ring}(R,+,\cdot)
\land
\operatorname{IdentityElement}(1,\cdot,R).
\]
\end{remark*}

\begin{remark*}[Definition predicate reading]
A ring has unity exactly when multiplication has a two-sided identity.
\end{remark*}

\begin{remark*}[Negated quantified statement]
\[
\neg\operatorname{RingWithUnity}(R,+,\cdot,1)
\Longleftrightarrow
\neg\operatorname{Ring}(R,+,\cdot)
\lor
\neg\operatorname{IdentityElement}(1,\cdot,R).
\]
\end{remark*}

\begin{remark*}[Interpretation]
The unity element is the multiplicative identity of the ring.
\end{remark*}

\begin{remark*}[Examples]
The integers, residue rings $\mathbb{Z}/n\mathbb{Z}$, polynomial rings, and
matrix rings are standard examples of rings.
\end{remark*}

\begin{dependencies}
\begin{itemize}
  \item \hyperref[def:algebraic-ring]{Ring}
  \item \hyperref[def:identity-element]{Identity Element}
\end{itemize}
\end{dependencies}
